插入排序（Insertion Sort）是最直观的排序算法之一。
它的工作方式很像你一手拿一叠扑克牌，另一只手把每一张牌插入到合适的位置。

\subsection{核心思路}

从第二个元素开始遍历数组。对于每个元素 \texttt{A[j]}：

\begin{enumerate}
    \item 把它暂存为 \texttt{key}
    \item 把 \texttt{key} 左边所有比它大的元素整体右移一格
    \item 最后把 \texttt{key} 放到腾出来的位置上
\end{enumerate}

这样，每处理完一个元素，它左边的部分就是有序的。

\subsection{时间复杂度}

\begin{center}
\begin{tabular}{lll}
最好情况 & 数组已经有序 & \(O(n)\) \\
最坏情况 & 数组完全逆序 & \(O(n^2)\) \\
平均情况 & 随机数据 & \(O(n^2)\) \\
\end{tabular}
\end{center}

对小规模数据或几乎有序的数据来说，插入排序实际上很快，
因为它空间开销为 \(O(1)\) 且常数因子很小。

\subsection{C 代码}

\lstinputlisting[style=cstyle, caption=src/sort/inserton\_sort.c — 插入排序]
{../src/sort/inserton_sort.c}

\subsection{代码解读}

这段代码的骨架非常简洁：

\begin{lstlisting}[style=cstyle]
void insertion_sort(int A[], int n) {
    for (int j = 1; j < n; j++) {      // 从第二个元素开始
        int key = A[j];                 // 拿出当前元素
        int i = j - 1;                  // i 指向已排序部分末尾
        while (i >= 0 && A[i] > key) {  // 向左找插入点
            A[i + 1] = A[i];            // 元素右移腾出空间
            i = i - 1;
        }
        A[i + 1] = key;                 // 放入正确位置
    }
}
\end{lstlisting}

几个值得注意的点：

\begin{itemize}
    \item \texttt{A[i] > key} 决定了排序是升序的；如果改成 \texttt{A[i] < key}，就变成降序。
    \item 内层 \texttt{while} 循环的「边比较边后移」是典型的原地插入方式，不需要额外空间。
    \item 整个算法是\textbf{稳定}的：相等元素不会互相交换相对位置。
\end{itemize}

编译运行：

\begin{lstlisting}[style=bashstyle]
clang -o insert src/sort/inserton_sort.c
./insert
2 4 5 6 1 3 7 8
\end{lstlisting}
